



\subsection{Single-Noise-Level Identification of the Loss Exponent at Finite Noise}


%We consider an interval stimulus space $\mathcal{X}=[0,1]$ with uniform encoding
%\begin{equation}
%F(\theta)=\theta,
%\end{equation}
%and an arbitrary smooth strictly positive asymmetric prior. For concreteness, take the exponential family
%\begin{equation}
%\Prior(\theta)=\frac{\lambda e^{\lambda \theta}}{e^\lambda-1}
%\end{equation}
%for some fixed $\lambda\neq 0$.
%Sensory encoding is
%\begin{equation}
%m=\theta+\varepsilon, \qquad \varepsilon\sim N(0,t),
%\end{equation}
%where $t>0$ is the sensory-noise variance.
%
%For uniform encoding, the single-noise low-noise confounding transformation from Theorem~\ref{thm:identifiability} becomes trivial, because
%\begin{equation}
%p_q(\theta)\propto \Prior(\theta) F'(\theta)^{\frac{\tilde q-\tilde p}{2}}=\Prior(\theta)
%\end{equation}
%for every pair of exponents $p,q$.
%Thus the infinitesimal-noise argument does not distinguish exponents here.
%However, the \emph{exact} finite-noise posterior does:
%\begin{align}
%p(\theta|m) &\propto \exp\left(-\frac{(\theta-m)^2}{2t}\right)e^{\lambda \theta}\mathbf{1}_{[0,1]}(\theta) \\
%&\propto \exp\left(-\frac{(\theta-\mu)^2}{2t}\right)\mathbf{1}_{[0,1]}(\theta),
%\end{align}
%where
%\begin{equation}
%\mu:=m+\lambda t.
%\end{equation}
%So the posterior is exactly a Gaussian $N(\mu,t)$ truncated to the interval $[0,1]$.
%
%At a fixed nonzero noise level $t$, different $L_p$ estimators select different points of this asymmetric truncated posterior.
%In particular:
%\begin{align}
%\hat\theta_{0}(m)&=\Pi_{[0,1]}(\mu), \\
%\hat\theta_{1}(m)&=\mu+\sqrt{t}\,\Phi^{-1}\left(\frac{\Phi(a)+\Phi(b)}{2}\right), \\
%\hat\theta_{2}(m)&=\mu+\sqrt{t}\,\frac{\phi(a)-\phi(b)}{\Phi(b)-\Phi(a)},
%\end{align}
%where
%\begin{equation}
%a=\frac{-\mu}{\sqrt{t}}, \qquad b=\frac{1-\mu}{\sqrt{t}},
%\end{equation}
%$\Pi_{[0,1]}$ denotes projection onto $[0,1]$, and $\phi,\Phi$ are the standard normal density and CDF.
%
%Now choose a measurement $m$ such that $\mu=1$, i.e.
%\begin{equation}
%m=1-\lambda t.
%\end{equation}
%Then
%\begin{align}
%\hat\theta_0 &= 1, \\
%\hat\theta_1 &= 1+\sqrt{t}\,\Phi^{-1}(1/4) = 1-0.67449\sqrt{t}, \\
%\hat\theta_2 &= 1-\sqrt{t}\,\frac{\phi(0)}{\Phi(0)} = 1-0.79788\sqrt{t}.
%\end{align}
%Therefore
%\begin{equation}
%\hat\theta_2 < \hat\theta_1 < \hat\theta_0.
%\end{equation}
%So the three exponents $p=2,1,0$ produce different deterministic decoders already at one fixed sensory-noise level.
%Because the full response distribution at that noise level determines the decoder, the exponent is identifiable from one level of sensory noise in this example.
%
%This example also clarifies the role of noise magnitude.
%As $t\to 0$, the separation between the estimators vanishes, recovering the low-noise confounding result.
%But for finite $t$, the gap between estimators is of order $\sqrt{t}$ near the boundary.
%Thus, departing from the infinitesimal-noise regime can \emph{improve} single-level identifiability of the loss exponent.
%This does not imply that identifiability improves monotonically for arbitrarily large noise; rather, it shows analytically that intermediate nonzero noise can reveal the exponent even when the infinitesimal-noise approximation cannot.
%
%\paragraph*{Circular Analogue}


Here we give an explicit analytical example showing that a single level of sensory noise can sometimes be sufficient for identifying the model.

The same conceptual mechanism also appears on a circular stimulus space, though the formulas are somewhat less explicit.
Take $\mathcal{X}=\mathbb{T}$, uniform encoding $F(\theta)=\theta$, and von Mises sensory noise
\begin{equation}
p(m|\theta)\propto \exp(\kappa \cos(\theta-m)),
\end{equation}
where $\kappa$ is the concentration parameter.
Choose the smooth strictly positive prior
\begin{equation}
\Prior(\theta)=\frac{1+\varepsilon \sin(2\theta)}{2\pi},
\qquad |\varepsilon|<1.
\end{equation}
This prior is periodic and similar to the prior used in (TODO XXX), but smooth and analytically more explicitly tractable.

Again the low-noise confounding transformation is trivial because $F'(\theta)\equiv 1$.
However, at finite $\kappa$, the exact posterior for a fixed measurement $m=0$ is
\begin{equation}
p(\theta|m=0)\propto e^{\kappa \cos\theta}\left(1+\varepsilon\sin(2\theta)\right),
\end{equation}
which is asymmetric as soon as $\varepsilon\neq 0$.

For small $\varepsilon$, one can compute the first-order displacement of different circular estimators explicitly.
If $\hat\theta_0$ denotes the MAP estimator, then differentiating the log posterior gives
\begin{equation}
\hat\theta_0=\frac{2\varepsilon}{\kappa}+O(\varepsilon^2).
\end{equation}
If $\hat\theta_2$ denotes the circular posterior mean direction, then
\begin{equation}
\hat\theta_2=
\arg \int e^{i\theta} e^{\kappa\cos\theta}\left(1+\varepsilon\sin(2\theta)\right)d\theta
=
\frac{\varepsilon}{2}\frac{I_1(\kappa)-I_3(\kappa)}{I_1(\kappa)}+O(\varepsilon^2),
\end{equation}
where $I_j$ are modified Bessel functions of the first kind.
Using the identity $I_1(\kappa)-I_3(\kappa)=\frac{4}{\kappa}I_2(\kappa)$, this becomes
\begin{equation}
\hat\theta_2=
\frac{2\varepsilon\, I_2(\kappa)}{\kappa I_1(\kappa)}+O(\varepsilon^2).
\end{equation}
Therefore
\begin{equation}
\hat\theta_0-\hat\theta_2=
\frac{2\varepsilon}{\kappa}\left(1-\frac{I_2(\kappa)}{I_1(\kappa)}\right)+O(\varepsilon^2),
\end{equation}
which is generically nonzero at any finite $\kappa$.
So already at one fixed noise level, the circular MAP and circular posterior mean differ.

Note that, by Theorem 1, the encoding and in particular $\kappa$, can be identified from a single noise level.

The distinction between the two estimators disappears as $\kappa\to 0$ (very large noise), because the posterior becomes nearly uniform and the two estimators nearly coincide with gap $\sim \varepsilon/2$.
The distinction between the two estimators disappears as $\kappa\to\infty$ (very small noise), because the two estimators nearly coincide with gap $\sim 3\varepsilon/\kappa^2$.

This reveals that identifiability of the loss exponent at a single noise level may be present at certain noise levels, but is not generically guaranteed at all noise levels.
Simulations show that there often is some noise level at which identifiability holds, but this noise level depends on the model.
For instance, Figure~\ref{fig:known-loss-unknown-prior} shows one model (uniform prior and periodic encoding) where identifiability holds at a single level of \emph{low} noise (top of that figure), and a model (periodic prior and uniform encoding) where a single level of \emph{high} noise is more successful (mid of that figure); in both cases, data from four noise levels is however much more successful at narrowing down the exponent closely (bottom of that figure).

%This circular example makes the same point as the interval example, but by a different mechanism.
%On the interval, finite-noise identifiability comes from boundary asymmetry of a truncated posterior.
%On the circle, there is no boundary, so the separation instead comes from prior-induced asymmetry of the posterior.
%In both cases, the distinction disappears in the infinitesimal-noise limit and becomes visible once one moves away from that limit.
